\section{材料の応力・変形}

\subsection{平成27年度 問6}
\label{q:h27-reg-applied-06}

\noindent\textbf{原典：} sources/H27/17\_applied\_science.pdf、PDF物理3頁。

\textbf{公式正答：} 5（single）。

\par\textit{原典本文の校正と独自解説は未作成。分類のみの登録項目。}

\section{電磁誘導・磁性}

\subsection{令和5年度 問7}
\label{q:r05-reg-applied-07}

\noindent\textbf{原典：} sources/R05/17\_applied\_science.pdf、PDF物理3頁。

\textbf{公式正答：} 4（single）。

\par\textit{原典本文の校正と独自解説は未作成。分類のみの登録項目。}

\section{原子核・放射線}

\subsection{平成23年度 問10}
\label{q:h23-reg-applied-10}

\noindent\textbf{原典：} sources/H23/17\_applied\_science.pdf、PDF物理4頁。

\textbf{公式正答：} 2（single）。

\par\textit{原典本文の校正と独自解説は未作成。分類のみの登録項目。}

\subsection{平成26年度 問9}
\label{q:h26-reg-applied-09}

\noindent\textbf{原典：} sources/H26/17\_applied\_science.pdf、PDF物理4頁。

\textbf{公式正答：} 2（single）。

\par\textit{原典本文の校正と独自解説は未作成。分類のみの登録項目。}

\subsection{平成28年度 問9}
\label{q:h28-reg-applied-09}

\noindent\textbf{原典：} sources/H28/17\_applied\_science.pdf、PDF物理4頁。

\textbf{公式正答：} 3（single）。

\par\textit{原典本文の校正と独自解説は未作成。分類のみの登録項目。}

\subsection{令和4年度 問9}
\label{q:r04-reg-applied-09}

\noindent\textbf{原典：} sources/R04/17\_applied\_science.pdf、PDF物理3頁。

\textbf{公式正答：} 2（single）。

\par\textit{原典本文の校正と独自解説は未作成。分類のみの登録項目。}

\section{光子の量子性・電子の物質波}

\subsection{平成24年度 問7}
\label{q:h24-reg-applied-07}

\noindent\textbf{原典：} sources/H24/17\_applied\_science.pdf、PDF物理2頁。

\textbf{公式正答：} 2（single）。

\paragraph{問題の要旨} 光の粒子性に直接関係しない現象又は概念を選ぶ。

\Needspace{4\baselineskip}\paragraph{確認した条件・記述（独自要約）}\mbox{}\par

\noindent\textbf{方向：} 光の粒子性に直接関係しないことがらとして最適。\par

\paragraph{問われる核心} 光子の粒子性と電子の物質波

\paragraph{必要知識} 光子の粒子性と電子の物質波。光子のエネルギーと運動量に関する1・3・4・5に対し、2は電子の波動性を観測する現象である。光の粒子性と物質粒子の波動性を区別して2を選ぶ。

\paragraph{考え方・正答を選ぶ理由（査読待ち初稿）} 光子のエネルギーと運動量に関する1・3・4・5に対し、2は電子の波動性を観測する現象である。光の粒子性と物質粒子の波動性を区別して2を選ぶ。

\Needspace{4\baselineskip}\paragraph{選択肢の検討（査読待ち初稿）}\mbox{}\par

\noindent\textbf{1：直接関係} 光電効果。光がエネルギーhf単位で電子へ渡すモデルで、閾値周波数と光電子の最大運動エネルギーを説明する。\par

\noindent\textbf{2：直接関係しない} 電子線回折。電子の物質波が結晶で干渉・回折する現象で、光の粒子性を直接示すものではない。\par

\noindent\textbf{3：直接関係} コンプトン効果。光子のエネルギー・運動量を電子との散乱で保存させ、波長変化を説明する。\par

\noindent\textbf{4：光電効果へ関係} 仕事関数。表面から電子を取り出す最小エネルギーで、光電効果のエネルギー収支へ入る。これ自体を光子と同一の概念にはしない。\par

\noindent\textbf{5：直接関係} プランク定数。光子エネルギーhfや運動量h/λを結ぶ定数で、光の量子性へ直接関係する。\par

\paragraph{再利用する考え方} 光子の粒子性と電子の物質波を暗記だけでなく、位相・幾何・符号と成立条件から導く。選択肢番号を類似年度へそのまま移さない。

\paragraph{共通理論} \hyperref[theory:modern:duality]{第\ref*{theory:modern:duality}節}

\par 参考：\cite{os_photoelectric,os_matter}。

\par 難易度 2（校準初稿）、優先度未評価。

\section{未確定}

\subsection{平成25年度 問7}
\label{q:h25-reg-applied-07}

\noindent\textbf{原典：} sources/H25/17\_applied\_science.pdf、PDF物理4頁。

\textbf{公式正答：} 5（single）。

\par\textit{原典本文の校正と独自解説は未作成。分類のみの登録項目。}

\subsection{平成27年度 問5}
\label{q:h27-reg-applied-05}

\noindent\textbf{原典：} sources/H27/17\_applied\_science.pdf、PDF物理2頁。

\textbf{公式正答：} 2（single）。

\par\textit{原典本文の校正と独自解説は未作成。分類のみの登録項目。}
